% Generated from the Lean dependency data; see the comment at the top of content.tex.
\chapter{Inner forms and the main theorem}
\label{chap:mainthm}
Thesis \S6.1--6.4. The Langlands parameter is not modelled: the main theorem is proved for every
family of packets satisfying the stability, transfer and genericity hypotheses, which is all the proof
uses.

\section{Inner forms, stability and packets}
\begin{definition}[Inner twists]
\label{InnerTwist}
\lean{GSpinGGP.InnerTwist}
\leanok
An inner twist of $\GSpin(W)$: a quadratic space $W'$ on $F^m$ of the same dimension and discriminant
with a fixed $\Fbar$-isometry $\varphi:W_{\Fbar}\simeq W'_{\Fbar}$ such that
$\varphi^{-1}\circ\sigma(\varphi)\in\SO(W)(\Fbar)$ for all $\sigma$. It gives
$\psi_\alpha=C(\varphi)\times C(\varphi\oplus\mathrm{id}_Z):G_{\Fbar}\simeq G_{\alpha,\Fbar}$.
This differs from Thesis \S6.2: the inner forms are indexed by $H^1(F,\SO(W))$, not $H^1(F,\mathrm{PGSpin}(W))$.
\end{definition}

\begin{definition}[Stable conjugacy across inner forms]
\label{StablyConjG}
\lean{GSpinGGP.GGPData.StablyConjG}
\leanok
\uses{G, GGPData, GSpin, InnerTwist}
$x\in G(F)$ and $y\in G_\alpha(F)$ are stably conjugate if $y=\psi_\alpha(g x g^{-1})$ for some
$g\in G(\Fbar)$, componentwise.
\end{definition}

\begin{definition}[The Kottwitz sign]
\label{kottwitzSign}
\lean{GSpinGGP.kottwitzSign}
\leanok
\uses{GGPData, InnerTwist}
$e(G_\alpha) = (-1)^{\mathrm{rk}_F G-\mathrm{rk}_F G_\alpha}=(-1)^{w(W)+w(W_\alpha)+w(V)+w(V_\alpha)}$
(the $p$-adic formula; that it equals Kottwitz's cohomological sign is \cite{Kot83}).
\end{definition}

\begin{definition}[Systems of inner forms]
\label{InnerFormSystem}
\lean{GSpinGGP.InnerFormSystem}
\leanok
\uses{GGPData, InnerTwist}
An index set $\{\alpha\}$ with inner twists $(W_\alpha,\varphi_\alpha)$, $\alpha=1$ the trivial one, such that every
nondegenerate form on $F^m$ with the discriminant of $W$ is isometric to exactly one $W_\alpha$; so the index
set is $H^1(F,\SO(W))=H^1(F,H/Z_{G,H})$. $(G_\alpha,H_\alpha,\xi_\alpha)$ is the GGP triple of
$(W_\alpha,V_\alpha=W_\alpha\oplus Z)$. This differs from Thesis \S6.2: the inner forms are indexed by
$H^1(F,\SO(W))$, not $H^1(F,\mathrm{PGSpin}(W))$.
\end{definition}

\begin{definition}[Stable functions]
\label{IsStableFn}
\lean{GSpinGGP.IsStableFn}
\leanok
\uses{G, GGPData, GSpin, InnerTwist, gspin}
(STAB) $\theta$ is constant on stable conjugacy classes of regular semisimple elements (Thesis \S6.1).
\end{definition}

\begin{definition}[Tempered packets on the inner forms]
\label{TemperedPacketFamily}
\lean{GSpinGGP.TemperedPacketFamily}
\leanok
\uses{G, GGPData, GSpin, InnerFormSystem, InnerTwist, IsStableFn, IsTempered, StablyConjG, UnitaryRep, WhittakerDatum, charG, gspin, kottwitzSign}
Finite families $\Pi_\alpha$ of pairwise inequivalent irreducible tempered representations of $G_\alpha(F)$
with central characters trivial on $Z_{G,H}(F)$, such that, with $\theta_\alpha=\sum_{\pi\in\Pi_\alpha}\theta_\pi$
(Thesis \S6.3):
(STAB) each $\theta_\alpha$ is stable;
(TRANS) $\theta_\alpha(y)=e(G_\alpha)\theta_1(x)$ for stably conjugate regular $x\in G(F)$, $y\in G_\alpha(F)$;
(WHITT) for every Whittaker datum of $G$ exactly one member of $\Pi_1$ has a Whittaker model.
\end{definition}

\section{The main theorem}
\begin{definition}[The multiplicity formula (statement)]
\label{MultiplicityFormula}
\lean{GSpinGGP.MultiplicityFormula}
\leanok
\uses{G, GGPData, GSpin, IsTempered, UnitaryRep, charG, mGeom, mult}
The statement of Thesis Thm.~7.2.2 (= Thm.~1.0.2): for every GGP triple and every irreducible tempered
$\pi$ with central character trivial on $A_{G,H}(F)$, $m(\pi)<\infty$ and $m(\pi)=\mgeom(\theta_\pi)$.
\end{definition}

\begin{theorem}[The multiplicity formula]
\label{multiplicity_formula}
\lean{GSpinGGP.multiplicity_formula}
\leanok
\uses{MultiplicityFormula}
The multiplicity formula holds (Thesis Thm.~7.2.2 = Thm.~1.0.2).
\end{theorem}
\begin{proof}
\leanok
\uses{multiplicityFormula_of_traceFormula}
By the strong induction on $\dim G$ of Thesis \S7.2, from the geometric expansion, the elliptic form of
Thesis Prop.~7.4.1(ii) and elliptic isolation (Chapter~\ref{chap:multformula}).
\end{proof}

\begin{theorem}[The count (Thesis (6.7.1))]
\label{sum_mult_eq_one}
\lean{GSpinGGP.sum_mult_eq_one}
\leanok
\uses{GGPData, InnerFormSystem, InnerTwist, IsQuasiSplit, TemperedPacketFamily, mult}
Let $G$ be quasi-split and $(\Pi_\alpha)$ as above, with finitely many inner forms. Then
$\sum_\alpha\sum_{\pi\in\Pi_\alpha} m(\pi)=1$.
\end{theorem}
\begin{proof}
\leanok
\uses{G, GSpin, IsTempered, LIT-germ_unique, LIT-hc_character_isQuasiCharacter, LIT-rodier_germ_one, UnitaryRep, WhittakerDatum, charG, germ, gspin, jHat, mGeom, multiplicity_formula, orbitChoice, orbitChoice_one_mem_nilReg, sum_mult_eq_germ_one}
By (6.7.7) the sum is the germ $c_{\theta_1,\mathcal O_1}(1)$. By Rodier's theorem applied to
the regular orbit $\mathcal O_1$, $c_{\theta_1,\mathcal O_1}(1)$ is the number of members of $\Pi_1$ with a Whittaker
model for the corresponding datum, which is $1$ by (WHITT). The multiplicity formula gives $m(\pi)<\infty$.
\end{proof}

\begin{theorem}[Main theorem: the local GGP conjecture for $\GSpin$]
\label{local_ggp}
\lean{GSpinGGP.local_ggp}
\leanok
\uses{GGPData, InnerFormSystem, InnerTwist, IsQuasiSplit, TemperedPacketFamily, mult}
(Thesis Thm.~6.4.1 = Thm.~1.0.1, $p$-adic case.) Let $G$ be quasi-split and $(\Pi_\alpha)_\alpha$ tempered
packets on the inner forms $G_\alpha$, $\alpha\in H^1(F,\SO(W))$, with central characters trivial on
$Z_{G,H}(F)$, satisfying (STAB), (TRANS), (WHITT). Then exactly one $\pi\in\bigsqcup_\alpha\Pi_\alpha$ has
$m(\pi)=1$, and $m(\pi')=0$ for all others.
\end{theorem}
\begin{proof}
\leanok
\uses{padic_quadratic_classes, sum_mult_eq_one}
The inner forms are finite in number (classification of $p$-adic quadratic forms). The sum of the
multiplicities is $1$, and each is a natural number, so exactly one is $1$ and the rest vanish
(Thesis \S6.7). Multiplicity one is not needed for this deduction.
\end{proof}

